\begingroup
\renewcommand{\arraystretch}{1.45}
\Needspace{46\baselineskip}
\InterpolationArticle{41}

Da alle in den Artikeln 29---40 betrachteten besonderen Fälle vollständig auf den allgemeinen Fall in den Artikeln 20---28 zurückgeführt sind, werden auch alle Kunstgriffe, durch die in diesem Fall die Rechnung abgekürzt wird, wie die von uns in den Artikeln 25, 26, 27 erläuterten, auf jene anwendbar sein. Daher wird es nicht nötig sein, bei dieser Untersuchung zu verweilen, die das folgende, zu den Artikeln 39, 40 gehörige Beispiel abschließen wird.

Die Mittelpunktsgleichung für den neuen Planeten \textit{Juno} ist mit der Exzentrizität \(0{,}254236\) nach der indirekten Methode in Schritten von jeweils zehn Grad wie folgt berechnet worden

\[
\begin{array}{r|r|rrrcr}
\multicolumn{2}{c|}{}&
\multicolumn{5}{c}{\text{Mittelpunktsgleichung}=X}\\
\multicolumn{2}{r|}{\text{Mittlere Anomalie}=x}&-&+&&&\\ \hline
0^{\mathord{0}}&360^{\mathord{0}}&0&&&&0\\
10&350&3^{\mathord{0}}&50'&38''30&\cdot\ \cdot\ \cdot\ \cdot&13838''30\\
20&340&7&38&21{,}47&\cdot\ \cdot\ \cdot\ \cdot&27501{,}47\\
30&330&11&20&8{,}79&\cdot\ \cdot\ \cdot\ \cdot&40808{,}79\\
40&320&14&52&48{,}06&\cdot\ \cdot\ \cdot\ \cdot&53568{,}06\\
50&310&18&12&49{,}21&\cdot\ \cdot\ \cdot\ \cdot&65569{,}21\\
60&300&21&16&17{,}02&\cdot\ \cdot\ \cdot\ \cdot&76577{,}02\\
70&290&23&58&42{,}92&\cdot\ \cdot\ \cdot\ \cdot&86322{,}92\\
80&280&26&14&55{,}85&\cdot\ \cdot\ \cdot\ \cdot&94495{,}85\\
90&270&27&58&52{,}36&\cdot\ \cdot\ \cdot\ \cdot&100732{,}36\\
100&260&29&3&28{,}13&\cdot\ \cdot\ \cdot\ \cdot&104608{,}13\\
110&250&29&20&33{,}68&\cdot\ \cdot\ \cdot\ \cdot&105633{,}68\\
120&240&28&41&2{,}10&\cdot\ \cdot\ \cdot\ \cdot&103262{,}10\\
130&230&26&55&22{,}77&\cdot\ \cdot\ \cdot\ \cdot&96922{,}77\\
140&220&23&55&2{,}70&\cdot\ \cdot\ \cdot\ \cdot&86102{,}70\\
150&210&19&35&0{,}79&\cdot\ \cdot\ \cdot\ \cdot&70500{,}79\\
160&200&13&57&40{,}52&\cdot\ \cdot\ \cdot\ \cdot&50260{,}52\\
170&190&7&16&58{,}33&\cdot\ \cdot\ \cdot\ \cdot&26218{,}33\\
180&180&0&&&&0
\end{array}
\]

\Needspace{15\baselineskip}
Wir zerlegen diese Periode von \(36\) Gliedern in sechs kleinere zu je sechs Gliedern; die sechs in jeder Periode enthaltenen Werte der Funktion \(X\) werden durch eine Formel von folgender Art dargestellt werden
\[
\begin{aligned}
&\gamma+\gamma'\cos x+\gamma''\cos 2x+\gamma'''\cos 3x\\[6pt]
&\quad+\delta'\sin x+\delta''\sin 2x+\delta'''\sin 3x
\end{aligned}
\]
\SourceLeaf{335}
wobei wir für die Koeffizienten \(\gamma,\gamma',\delta'\) usw. die folgenden Werte finden:


\begin{landscape}\thispagestyle{empty}\begingroup\fontsize{9}{12}\selectfont\setlength{\arraycolsep}{3pt}\renewcommand{\arraystretch}{1.3}
\[
\begin{array}{l|r|r|r|r|r|r|r|r}
\text{Periode}&y=6x&\gamma&\gamma'&\gamma''&\gamma'''&\delta'&\delta''&\delta'''\\ \hline
\text{erste}&0^{\mathord{0}}&0&0&0&0&
-103830{,}165&
+15406{,}638&0\\
\text{zweite}&60&+56''205&-217''757&+761''671&-1528{,}825&-103937{,}346&+15841{,}387&-882{,}667\\
\text{dritte}&120&+56{,}115&-217{,}467&+760{,}805&-1527{,}397&-104151{,}314&+16709{,}402&-2645{,}530\\
\text{vierte}&180&0&0&0&0&-104258{,}100&+17142{,}684&-3525{,}740\\
\text{fünfte}&240&-56{,}115&-217{,}467&-760{,}805&+1527{,}397&-104151{,}314&+16709{,}402&-2645{,}530\\
\text{sechste}&300&-56{,}205&-217{,}757&-761{,}671&+1528{,}825&-103937{,}346&+15841{,}387&-882{,}667
\end{array}
\]
\endgroup\end{landscape}


\Needspace{39\baselineskip}
Die einzelnen Koeffizienten \(\gamma,\gamma',\gamma'',\gamma''',\delta',\delta'',\delta'''\) werden wiederum auf eine Form von folgender Art zurückgeführt werden
\[
\begin{aligned}
&\varepsilon+\varepsilon'\cos 6x+\varepsilon''\cos 12x+\varepsilon'''\cos 18x\\[6pt]
&\quad+\zeta'\sin 6x+\zeta''\sin 12x+\zeta'''\sin 18x
\end{aligned}
\]
es lässt sich aber ohne Mühe erkennen, dass für die ersten vier \(\varepsilon,\varepsilon',\varepsilon'',\varepsilon'''\) verschwinden müssen; und für die letzten drei \(\zeta',\zeta'',\zeta'''\). Auf diese Weise findet sich
\[
\begin{aligned}
&\gamma=+64''{,}848\sin 6x+0''{,}052\sin 12x\\[6pt]
&\gamma'=-251''{,}277\sin 6x-0''{,}167\sin 12x\\[6pt]
&\gamma''=+879''{,}002\sin 6x+0''{,}500\sin 12x\\[6pt]
&\gamma'''=-1764''{,}511\sin 6x-0''{,}824\sin 12x\\[6pt]
&\delta'=-104044''{,}264+213''{,}968\cos 6x+0''{,}132\cos 12x\\[6pt]
&\qquad{}+0''{,}000\cos 18x\\[6pt]
&\delta''=+16275''{,}150-868''{,}020\cos 6x-0''{,}489\cos 12x\\[6pt]
&\qquad{}-0''{,}003\cos 18x\\[6pt]
&\delta'''=-1763''{,}689+1762''{,}868\cos 6x+0''{,}819\cos 12x\\[6pt]
&\qquad{}+0''{,}002\cos 18x
\end{aligned}
\]

\Needspace{24\baselineskip}
Setzt man diese Werte für \(\gamma,\gamma'\) usw. ein und beachtet die Regeln des vorangehenden Artikels, so ergibt sich für die Mittelpunktsgleichung die folgende Funktion, in der die einzelnen Koeffizienten bis auf ein Hundertstel einer Bogensekunde genau sind.


\[
\begin{array}{r|r}
-104044''264\sin x&-1''643\sin 9x\\
+16275{,}150\sin 2x&+0{,}494\sin 10x\\
-3527{,}378\sin 3x&-0{,}149\sin 11x\\
+873{,}511\sin 4x&+0{,}052\sin 12x\\
-232{,}622\sin 5x&-0{,}017\sin 13x\\
+64{,}848\sin 6x&+0{,}006\sin 14x\\
-18{,}655\sin 7x&-0{,}004\sin 15x\\
+5{,}491\sin 8x&+0{,}003\sin 16x
\end{array}
\]


\Needspace{36\baselineskip}
\SourceLeaf{336}
Auf diese Weise verhält sich die Rechnung offenbar so, als würde anstelle des Bogens \(x\) ein anderer, \(x'=x-a\), eingeführt. Häufig kann die Rechnung auch sehr erleichtert werden, wenn anstelle der Funktion \(X\) eine andere, \(X-\mathfrak{x}\), verwendet wird, wobei für \(\mathfrak{x}\) jede beliebige der Funktion \(X\) ähnliche Funktion angenommen werden kann, sofern sie die \(\frac12\mu-\frac12{}^{\mathrm{te}}\) Ordnung bei einem ungeraden Wert oder die \(\frac12\mu^{\mathrm{te}}\) Ordnung bei einem geraden Wert von \(\mu\) nicht überschreitet und im letzteren Fall die Koeffizienten der Glieder \(\cos\frac12\mu x\) und \(\sin\frac12\mu x\) (falls sie so weit ansteigt) im erforderlichen Verhältnis \((\cos\frac12\mu a:\sin\frac12\mu a)\) stehen. Offenbar werden, wenn die Werte der Funktion \(\mathfrak{x}\) für \(x=a.b,c,d\) usw. der Reihe nach \(\mathfrak{A},\mathfrak{B},\mathfrak{C},\mathfrak{D}\) usw. sind, die Werte der Funktion \(X-\mathfrak{x}\) für dieselben Werte der Reihe nach \(A-\mathfrak{A},B-\mathfrak{B},C-\mathfrak{C},D-\mathfrak{D}\) usw. sein; wird ferner die diesen Werten genügende Funktion \(Z\) nach den vorangehenden Regeln aufgesucht, so wird die Funktion \(\mathfrak{x}+Z\) mit derjenigen identisch sein, die aus der unmittelbaren Anwendung dieser Regeln auf die Werte der Funktion \(X\) selbst hervorginge. Die Funktion \(\mathfrak{x}\) muss aber stets so bestimmt werden, dass die vorgelegten Werte der Funktion \(X\) durch sie nahezu wiedergegeben werden, wodurch die Größen \(A-\mathfrak{A},B-\mathfrak{B},C-\mathfrak{C},D-\mathfrak{D}\) usw. klein und für die Rechnung leichter handhabbar werden. Übrigens kann für \(\mathfrak{x}\) auch eine konstante Größe angenommen werden, z. B. diese, \(\frac1\mu(A+B+C+D+\text{usw.})\), die den unveränderlichen Teil der gesuchten Funktion bildet.

\begin{center}\rule{0.28\linewidth}{0.4pt}\end{center}
\endgroup
